\subsection{対角基礎代数の原包絡への核は二通りだけである}
\label{subsec:h3-positive3-base-kernel}
\begin{lemma}\label{lem:h3-qweyl-domain-quotients}
$A=k\langle u,w\rangle/(wu-quw-c)$、$c\ne0$とする。
$q=1$なら$A$は単純である。
$q\ne1$では$d=c+(q-1)uw$と置くと、
任意の非零完全素イデアル$J$は$d$を含む。
また、パラメータが1でない二変数量子torusは非零完全素イデアルを持たない。
ここで完全素とは商が整域であることを意味する。
\end{lemma}
\begin{proof}
$q=1$の単純性は、PBW多項式を$u,w$との交換子で順に微分して
任意の非零イデアルに非零定数を作る通常のWeyl環の議論による。
$q\ne1$では$du=qud,dw=q^{-1}wd$であり、$d$をOre局所化できる。

$q$が無限位数なら、
\[
 A[d^{-1}]=k[d,d^{-1}](\sigma,t),\quad
 \sigma(d)=qd,\quad t=(d-c)/(q-1),\quad X=w,Y=u
\]
という一次GWA表示を持つ。
$k[d^{\pm1}]$には非零固有$\sigma$不変イデアルがなく、
$\sigma$の正冪は恒等でない。
$t,\sigma^nt$の零点も異なるので全$n>0$で互いに素である。
Bavulaの単純性定理\cite[Theorem 2.1]{BavulaGWAOverview}の四条件が全て満たされ、
$A[d^{-1}]$は単純となる。従って非零完全素$J$は$d$を含む。

$q$が位数$N>1$の場合、PBWの最高項から
\[
 Z(A)=k[U,W],\qquad U=u^N,W=w^N
\]
で、$A$はこの中心上自由階数$N^2$である。
\[
 d^N=c^N+(-1)^{N-1}(q-1)^NUW
\]
は$u^Nw^N=\prod_{j=0}^{N-1}(q^{-j}d-c)/(q-1)^N$から従う。
$A[d^{-1}]$は中心$R=k[U,W,(d^N)^{-1}]$上Azumaya階数$N^2$である。
実際、閉中心点で$U\ne0$なら、$d$の相異なる$N$固有値と巡回可逆$u$を取り、
$w=u^{-1}(d-c)/(q-1)$とすれば全行列環を得る。
$W\ne0$も対称的である。
$U=W=0$では
\[
 ue_j=e_{j+1},\qquad
 we_j=c(1+q+\cdots+q^{j-1})e_{j-1},\quad0\le j<N
\]
を端で零とすればよい。全ての内部係数が非零なので、これも全行列環となる。
各fiberの次元は$N^2$なのでこれらの表示は同型であり、Azumaya性が従う。

もし非零完全素$J$が$d$を含まなければ、局所化後も非零で整域商を持つ。
Azumaya代数の両側イデアルと中心イデアルの対応により、
非零素イデアル$\mathfrak p\subset R$上のAzumaya商となる。
$K=\operatorname{Frac}(R/\mathfrak p)$は超越次数1以下で、
この商のgeneric fiberは次数$N$の中心的単純代数である。
整域商を局所化した有限次元代数なのでdivisionであるはずだが、
Tsenの定理による$\operatorname{Br}(K)=0$
\cite[Tag 03RF]{StacksTsenBrauer}により$M_N(K)$へ分裂し矛盾する。
超越次数0では$K=k$なので同じ結論である。
従って$d\in J$である。

量子torusのパラメータが無限位数なら、最少単項式数の元に交換子を施す通常の議論で単純。
有限位数$M>1$なら、中心Laurent二変数環上Azumaya次数$M$である。
閉fiberの同定はclock/shift行列により直接得られる。
非零完全素商があれば中心の超越次数が1以下となり、同じTsenの議論に反する。
\end{proof}

\begin{theorem}\label{thm:h3-positive3-base-kernel-dichotomy}
命題\ref{prop:h3-positive3-nonzero-gamma-base}の$S$から整域への次数付き準同型を考え、
$z$の像は非零とする。その核$I$は
\[
 I=0\qquad\text{または}\qquad I=(C'),\quad C'=xy-a^{-2}yx
\]
に限られる。$a^4=1$では必ず$I=0$である。
$a^4\ne1$の第二の場合、像は二変数量子平面の偶Veronese環である。
従ってH3包絡への対角基礎写像にも、これ以外の高次核は存在しない。
これは追加の正生成元$v$の関係やH3許容条件を分類したという意味ではない。
\end{theorem}
\begin{proof}
$z$は$S$で正規であり、$I$は完全素なので$z$で飽和している。
$u=xz^{-1},w=yz^{-1}$とすると、順逆の生成元対応により
\[
 S[z^{-1}]=A_q(c)[z,z^{-1};\theta],\quad
 q=a^{-4},\quad c=\gamma a^{-3},\quad
 \theta(u,w)=(au,a^{-1}w)
\]
である。ここで$A_q(c)$は前補題の量子Weyl環である。
$z$の次数が1、$u,w$の次数が0なので、この局所化はstrongly $\mathbb Z$-graded。
従って$I\ne0$なら次数零部分の核$J\subset A_q(c)$も非零で、完全素である。
$a^4=1$なら$q=1$で量子Weyl環は単純なので、このような$I$はない。

$a^4\ne1$では前補題から$d\in J$である。
\[
 d=c+(q-1)uw=-aC'z^{-2}
\]
なので$C'\in I$。
$A_q(c)/(d)=k[u,u^{-1}]$であり、$z$を加えた商は
$zu=auz$という二変数量子torusである。
前補題によりこれには非零完全素イデアルがないので、
$I$の局所化は正確に$(C')$の局所化である。

最後に$S/(C')$を具体的に確認する。
$\kappa=(a^4-1)/\gamma\ne0$とすると、その関係は
\[
 yx=a^2xy,\quad zx=axz,\quad yz=azy,\quad z^2=\kappa xy.
\]
$x^iz^jy^k$、$j\in\{0,1\}$がPBW基底となる。
$\rho^2=a^{-1},\lambda^2=\rho\kappa$を取り、
$st=\rho ts$を満たす量子平面内で
\[
 x=s^2,\quad y=t^2,\quad z=\lambda st
\]
と置けば、これら全関係を満たす。
両者のHilbert級数は$(1+t)/(1-t)^2$なので、この対応は偶Veroneseとの同型である。
特に$S/(C')$は整域であり、$(C')$は$z$で飽和している。
局所化での等式を戻して$I=(C')$を得る。
\end{proof}

\begin{corollary}\label{cor:h3-positive3-base-survives-cut}
H3の$\gamma\ne0$対角chartでは、基礎代数$S$の像は三次切断
$B\to R=B/(G)$の後も単射的に保たれる。
特に$G$は$B$内の三生成元部分代数$\langle x,y,z\rangle$に属さない。
従って命題\ref{prop:h3-positive3-base-normal-cubics}の原始6乗根条件を
全H3に対するパラメータ制約として適用しない。
\end{corollary}
\begin{proof}
$B,R$は共に整域で、$z$の一次の像は非零である。
従って$S\to B$と$S\to R$の核は前定理の二通りだけである。
$G$は三次なので$B_2=R_2$であり、両核の二次部分は同じ。
二つの可能な核はともにその二次部分で決まるから、両核は完全に等しい。
従って$B$内の基礎像は$R$にも埋まり、非零$G$がその像に入ることはない。
\end{proof}

\texttt{verify\_cubic\_positive3\_base\_kernel.py}は
量子Weyl局所化、$d$正常化、GWA式と全軌道に用いるresultant、
根位数一般のoscillator漸化式、coneの四overlapとHilbert値を厳密検算する。
有限位数全体のAzumaya性とTsenの適用は上の構造証明による。
Bavulaの著者原稿Theorem 2.1とStacks Project Tags 03RD/03RFを一次資料で確認した。
